% ===================================================================
%  TABELLA N-RO 6 — Le interior del casa, le mobiliario, le
%  alimentation, le illumination.
%  Traduction 2026 d'apres texto/io, controlee sur texto/fr et
%  texto/es. Consigne : voir texto/ia/10-tabella-01.tex.
%
%  Deux notes, toutes deux marquees « (*) », et deux appels. L'ordre
%  les departage, comme dans les autres colonnes.
%
%  LA FEMME DE CHAMBRE PREND SON (16), contre le fac-simile ido qui
%  ecrit (6) — son propre numero du savon. C'est le francais qui a
%  raison, et l'ecart est declare dans outils/renvoji.py.
% ===================================================================

%%K t06-apar-1 apar
\VUsaut{6.0mm}
\VUtitre{15.0pt}{60}{TABELLA N\textsuperscript{o} 6}
\VUsaut{1.4mm}
\VUtitre{11.4pt}{0}{Le interior del casa, le mobiliario, le}
\VUsaut{0.4mm}
\VUtitre{11.4pt}{0}{alimentation, le illumination.}
\VUsaut{2.2mm}

%%K t06-apar-2 apar
Le casa es finite. Le oncle de Johannes es installate in su nove domicilio, del
qual nos cognosce le disposition. Ora videamus como ille mobilava su casa. Nos
explorara primo:

%%K t06-tit-1 sub
\VUpk{10.2pt}{40}{Le camera de dormir.}
\VUsaut{3.0mm}

%%K t06-c1-01-1 p
1. --- Nos arriva a un mal momento, proque le camerera es occupate a nettar le
\VUgras{camera}.

%%K t06-c1-02-1 p
2. --- Sur le \VUgras{tabula de toilette} \textsuperscript{(1)} il ha hic e
illac diverse objectos con le quales on se lava, se pectina, in un parola on
face su toilette. Illos es le \VUgras{bassinetto} \textsuperscript{(2)} e le
\VUgras{urceo de aqua} \textsuperscript{(3)}, le \VUgras{brossa pro capillos}
\textsuperscript{(4)}, le \VUgras{pectine} \textsuperscript{(5)}, le
\VUgras{sapon} \textsuperscript{(6)} e le \VUgras{spongia}
\textsuperscript{(7)}, finalmente un \VUgras{flacon} \textsuperscript{(8)} pro
le dentifricio liquide.

%%K t06-c1-03-1 p
3. --- Sur le solo, ante le tabula de toilette, ecce le \VUgras{situla}
\textsuperscript{(9)} pro colliger le aquas immunde.

%%K t06-c1-04-1 p
4. --- Nos vide in le \VUgras{antecamera} \textsuperscript{(10)} alcun
\VUgras{crocos pro vestimentos} \textsuperscript{(13)} e duo \VUgras{parapluvias}
\textsuperscript{(11)} in le \VUgras{portaparapluvias}
\textsuperscript{(12)}.

%%K t06-c1-05-1 p
5. --- Al altere latere del porta il ha le \VUgras{armario a speculo}
\textsuperscript{(14)} que contine le \VUgras{linage} \textsuperscript{(15)}.

%%K t06-c1-06-1 p
6. --- Profitemus que le \VUgras{camerera} \textsuperscript{(16)} arrangia le
\VUgras{lecto} \textsuperscript{(17)}, pro explorar su diverse partes. Le
\VUgras{quadro del lecto} \textsuperscript{(20)} contine un bon \VUgras{matras a
resortos} \textsuperscript{(21)} sur le qual Justina tosto ponera un
\VUgras{matras} \textsuperscript{(22)} de lana. Postea illa prendera del sedias
le duo \VUgras{drappos de lecto} \textsuperscript{(23)} e le \VUgras{coperilecto}
\textsuperscript{(24)}. Alora illa ponera sur le testiera, le \VUgras{cossino
transverse} \textsuperscript{(25)}, le \VUgras{cossino de capite}
\textsuperscript{(26)} que es con le \VUgras{plumon} \textsuperscript{(27)} sur
le \VUgras{tapete del lecto} \textsuperscript{(32)}. Illa usa un plumaceo pro
dispulverar le \VUgras{cortinas} \textsuperscript{(19)} e le \VUgras{baldachino
del lecto} \textsuperscript{(18)} e ecce un parte del labor finite.

%%K t06-c1-07-1 p
7. --- Alora illa debera ordinar un pauco le \VUgras{tabuletta de nocte}
\textsuperscript{(28)}, retender le \VUgras{eveliator} \textsuperscript{(29)},
preparar le \VUgras{parve lampa de nocte} \textsuperscript{(30)}, levar le
\VUgras{candela} \textsuperscript{(31)} pro nettar le candeliero.

%%K t06-tit-2 sub
\VUsaut{5.0mm}
\VUpk{10.2pt}{40}{Le corbe a labor e le accessorios del camino.}
\VUsaut{3.0mm}

%%K t06-c1-08-1 p
8. --- Le \VUgras{corbe a labor} \textsuperscript{(34)} presso le
\VUgras{tabouret} \textsuperscript{(33)} anque ha multe besonio de esser
ordinate. Illa debera plicar le \VUgras{brodaturas} \textsuperscript{(35)},
inserrar in le \VUgras{tabula de labor} \textsuperscript{(38)} le
\VUgras{digital}, le \VUgras{filo} \textsuperscript{(36)}, le \VUgras{acos}
\textsuperscript{(37)} e le \VUgras{cisorios} \textsuperscript{(39)} e clauder
con le \VUgras{clave} \textsuperscript{(40)} le coperculo del tabula.

%%K t06-c1-09-1 p
9. --- Sur le \VUgras{gueridon} \textsuperscript{(41)} del medio, Justina
renovara le aqua del \VUgras{carrafa} \textsuperscript{(42)}. Illa ponera
\VUgras{sucro} in le \VUgras{sucriera} \textsuperscript{(43)}, nettara le
\VUgras{pincetta pro sucro} \textsuperscript{(44)} e levara le \VUgras{brossa pro
vestimentos} \textsuperscript{(46)}.

%%K t06-c1-10-1 p
10. --- Le latere dextre del camera requirera de illa minus labor, proque le
\VUgras{chaise-longue} \textsuperscript{(47)} e su \VUgras{cossino}
\textsuperscript{(48)} es a lor juste loco, e, proque il non face frigido, nihil
es displaciate inter le accessorios (utensiles) del \VUgras{camino}
\textsuperscript{(49)}. Le \VUgras{pala} \textsuperscript{(50)}, le
\VUgras{tenalias} \textsuperscript{(51)}, le \VUgras{chenettos}
\textsuperscript{(55)}, le \VUgras{guarda-cinere} \textsuperscript{(56)} es
nette; le \VUgras{schermo de foco} \textsuperscript{(54)} non esseva displaciate
e le \VUgras{cassa de ligno} \textsuperscript{(52)} es ben providite de
\VUgras{ligno a arder} \textsuperscript{(53)}.

%%K t06-noto-1 noto
\VUnotes{143.2mm}{%
(*) Babo = infantetto, le A. baby, F. bébé, I. bambino.%
}

%%K t06-noto-2 noto
\VUnotes{143.2mm}{%
(*) Lambrequino = F. lambrequin, E. lambrequines, Port. lambrequins, mesmo G.
A. lambrequins, dunque G. A. F. P. E.%
}

%%K t06-c1-11-1 p
11. --- Tamen illa debera dispulverar le \VUgras{pendula}
\textsuperscript{(57)}, le \VUgras{candelabros} \textsuperscript{(58)} e le
\VUgras{portaobjectos} \textsuperscript{(59)}, finalmente ablue le
\VUgras{speculo} \textsuperscript{(60)}.

%%K t06-c1-12-1 p
12. --- Quanto al \VUgras{cuna} \textsuperscript{(61)} del infantetto (del babo
(*)), illo ja es preste.

%%K t06-tit-3 sub
\VUsaut{5.0mm}
\VUpk{10.2pt}{40}{Le salon.}
\VUsaut{3.0mm}

%%K t06-c2-01-1 p
1. --- Ora ecce le salon. Illo es un multo belle \VUgras{camera}. Le camino, que
se trova al angulo dextre es ornate de un delicate \VUgras{statuetta}
\textsuperscript{(1)}, de duo belle \VUgras{vasos} \textsuperscript{(3)} e
drappate de un \VUgras{lambrequino} \textsuperscript{(2)} de veluto (*).

%%K t06-c2-02-1 p
2. --- Pro impedir que le scintillas del focar arde le tapete, on poneva ante le
focar un \VUgras{para-scintillas} \textsuperscript{(4)} de cupro.

%%K t06-c2-03-1 p
3. --- Presso, un elegante \VUgras{scriptorio} \textsuperscript{(5)} de dama es
aperte. Sur illo le \VUgras{domina del casa} \textsuperscript{(23)} scribe su
litteras.

%%K t06-c2-04-1 p
4. --- Le fenestra es guarnite de \VUgras{cortinas de tulle}
\textsuperscript{(6)} con \VUgras{ligamines} \textsuperscript{(8)} appendite a
\VUgras{crocos} \textsuperscript{(7)} e de grande cortinas de stoffa
\VUgras{drappate in alto} \textsuperscript{(9)}.

%%K t06-c2-05-1 p
5. --- In le embrasura del fenestra, un \VUgras{cache-pot}
\textsuperscript{(11)} contine un \VUgras{planta verde} \textsuperscript{(10)},
un splendide palmiero del qual le folios se extende quasi usque al
\VUgras{lampa a petroleo} \textsuperscript{(12)}.

%%K t06-c2-06-1 p
6. --- Le \VUgras{umbrella del lampa} \textsuperscript{(13)} es arrangiate per
Maria, le incantatori \VUgras{juvena} \textsuperscript{(14)} qui es al piano
\textsuperscript{(15)}. Sedente multo recte sur le \VUgras{tabouret}
\textsuperscript{(16)}, illa studia un partitura. Illa es multo habile e curose,
como proba su \VUgras{armario a musica} \textsuperscript{(17)} que es
perfectemente ordinate.

%%K t06-tit-4 sub
\VUsaut{5.0mm}
\VUpk{10.2pt}{40}{Le mal puero.}
\VUsaut{3.0mm}

%%K t06-c2-07-1 p
7. --- Su fratre, Paulo, cerca un \VUgras{libro} \textsuperscript{(19)} in le
\VUgras{bibliotheca} \textsuperscript{(18)}. Ille es un \VUgras{juvene}
\textsuperscript{(20)}, movimentose e intolerabile. Suspendente se al bibliotheca
pro attinger le libro que es troppo alte, ille risca facer cader le
\VUgras{busto} \textsuperscript{(21)} que es supra.

%%K t06-c2-08-1 p
8. --- Ille profita, pro facer iste stultessa, que su matre es occupate a
conversar con un amica, un \VUgras{dama} \textsuperscript{(22)} qui veniva
passar alcun dies in su casa. Le matre sede sur un \VUgras{canape}
\textsuperscript{(24)} presso le \VUgras{fauteuil} \textsuperscript{(25)}, e su
amica es sur le \VUgras{sedia} \textsuperscript{(27)} del qual on vide solmente
le \VUgras{dorso} \textsuperscript{(26)}.

%%K t06-c2-09-1 p
9. --- Le grande \VUgras{quadro} \textsuperscript{(30)} pendente al muro contine
le \VUgras{portrait} \textsuperscript{(29)} de un parenta. In le
\VUgras{album} \textsuperscript{(32)} ponite sur le tabula es reunite le
\VUgras{photographias} \textsuperscript{(33)} del altere membros del familia. Le
infantes justo lo foliava, proque le \VUgras{sedias} \textsuperscript{(31)} es
ancora gruppate circa le tabula.

%%K t06-c2-10-1 p
10. --- Le \VUgras{vaso chinese} \textsuperscript{(34)} de porcellana, que es
presso le \VUgras{cultello a papiro} \textsuperscript{(35)} esseva donate a Maria
pro su festa, e le \VUgras{paravento} \textsuperscript{(36)} que guarni le angulo
del camera esseva apportate de China per su oncle.

%%K t06-c2-11-1 p
11. --- Allegrate per le belle \VUgras{picturas} \textsuperscript{(37)}
appendite al \VUgras{pariete} \textsuperscript{(42)}, illuminate per le
\VUgras{lustro del plafon} \textsuperscript{(38)}, del qual le \VUgras{lampas
electric} \textsuperscript{(39)} le vespere brilla gaimente, iste salon pare
multo agradabile. Le molle \VUgras{tapete} \textsuperscript{(40)} extendite sur
le parquete, e le tanto facile a usar \VUgras{campana electric}
\textsuperscript{(41)}, completa su conforto.

%%K t06-tit-5 sub
\VUsaut{3.6mm}
\VUpk{10.2pt}{40}{Le sala a mangiar.}
\VUsaut{3.0mm}

%%K t06-c3-01-1 p
1. --- Sonava le hora del cena. Tote le familia, excepte Maria, es reunite circa
le tabula. Le sala a mangiar es agradabilemente calefacite per le excellente
\VUgras{stufa} \textsuperscript{(1)} de faience, con un tenue \VUgras{tubo}
\textsuperscript{(2)}.

%%K t06-c3-02-1 p
2. --- Iste camera non contine multe mobiles. Al longo del muro pende, super le
\VUgras{banchetta} \textsuperscript{(4)}, un \VUgras{fontanetta}
\textsuperscript{(3)} ornamental. Multo presso, il ha le \VUgras{credentia}
\textsuperscript{(5)} sur le qual on pone le alimentos frigide, le plattos de
dessert, e le diverse utensiles de tabula de que on non ha besonio
immediatemente.

%%K t06-c3-03-1 p
3. --- Assi nos vide sur le tabuletta superior, un \VUgras{pullo rostite}
\textsuperscript{(7)}, un \VUgras{pisce} \textsuperscript{(8)} e un
\VUgras{cuppa} \textsuperscript{(10)} con \VUgras{fructos}
\textsuperscript{(9)}. Infra ecce le \VUgras{caseo} \textsuperscript{(11)}, le
\VUgras{pan} \textsuperscript{(12)}, le \VUgras{portaflacones}
\textsuperscript{(13)} que contine toto lo que on besonia pro condimentar le
salata: le oleo, le vinagre, le \VUgras{sal} \textsuperscript{(14)} e le
\VUgras{pipere} \textsuperscript{(15)}. Finalmente, sur le tabuletta inferior il
ha un \VUgras{torta} \textsuperscript{(17)} presso le \VUgras{mustardiera}
\textsuperscript{(16)}.

%%K t06-c3-04-1 p
4. --- Le horologio del sala a mangiar, que on nomina \VUgras{horologio de muro}
\textsuperscript{(20)}, ha un forma multo particular. Illo es incastrate in un
quadro de ligno que contine ancora un \VUgras{barometro}
\textsuperscript{(19)} e un \VUgras{thermometro} \textsuperscript{(18)}.

%%K t06-tit-6 sub
\VUsaut{4.6mm}
\VUpk{10.2pt}{40}{Le servicio del cena.}
\VUsaut{3.0mm}

%%K t06-c3-05-1 p
5. --- A tote le muros del camera es applicate \VUgras{pannellos de ligno}
\textsuperscript{(21)} de querco cerate. Un \VUgras{portiera}
\textsuperscript{(22)} de stoffa claude satis ben le porta que da entrata al
veranda. Illac le familia ira biber le caffe post le cena. Ja le \VUgras{tassas}
\textsuperscript{(24)} e le \VUgras{sub-tassas} \textsuperscript{(25)} es preste
sur le \VUgras{mobile a plattos} \textsuperscript{(23)}.

%%K t06-c3-06-1 p
6. --- Super le tabula es fixate al plafon un \VUgras{suspension}
\textsuperscript{(26)} del qual le multo brillante lampa le vespere illumina
tote le sala a mangiar.

%%K t06-c3-07-1 p
7. --- Ora exploremus le \VUgras{tabula a mangiar} \textsuperscript{(27)} mesme.
Quando le familia entrava, le coperto esseva preste. Sur le \VUgras{coperi-tabula}
\textsuperscript{(28)} esseva ponite, al loco de cata conviva, un
\VUgras{platto} \textsuperscript{(33)}, un \VUgras{servietta}
\textsuperscript{(29)}, un \VUgras{coclear} \textsuperscript{(34)}, un
\VUgras{furca} \textsuperscript{(35)}, un \VUgras{cultello}
\textsuperscript{(36)} e un \VUgras{vitro} \textsuperscript{(32)}. Il habeva
anque \VUgras{bottilias} \textsuperscript{(30)} pro le vino e un
\VUgras{carrafa} \textsuperscript{(31)} pro le aqua.

%%K t06-c3-08-1 p
8. --- Le \VUgras{servitor} \textsuperscript{(39)} primo apportava le
\VUgras{suppiera} \textsuperscript{(37)}, in le qual ancora fuma un pauco de
suppa. Ora ille servi le prime \VUgras{platto} \textsuperscript{(38)}, un platto
de carne. Postea ille presentara illos que nos ja nominava, finalmente, post le
\VUgras{platto de dessert}, ille profitara que su patronos habera entrate le
veranda, pro levar le utensiles de tabula e reordinar le sala a mangiar.

%%K t06-c3-08-2 p
\textit{Nota.} --- \textit{Le parolas de uso commun concernente le alimentos, es
hic indicate multo summarimente. Le lista essera completate sur le duo tabellas
« Le Hotel » n\textsuperscript{o} 13) e « Le Mercato »
(n\textsuperscript{o} 15).}

%%K t06-tit-7 sub
\VUsaut{4.6mm}
\VUpk{10.2pt}{40}{Le cocina.}
\VUsaut{3.0mm}

%%K t06-c4-01-1 p
1. --- In le cocina, que nos comencia a reguardar tra le porta medio aperte, nos
vide un pittoresc disordine. Ante le \VUgras{focar} \textsuperscript{(1)} in le
qual crepita un \VUgras{foco} \textsuperscript{(3)}, es installate un
\VUgras{rostitorio} \textsuperscript{(5)}. Un belle \VUgras{rostito}
\textsuperscript{(7)} es \VUgras{sur le spido} \textsuperscript{(6)} e fini
cocer se.

%%K t06-c4-02-1 p
2. --- Le \VUgras{sufflatorio} \textsuperscript{(8)} e le \VUgras{pala pro
carbon} \textsuperscript{(9)}, que on justo usava, restava sur le solo como le
\VUgras{ligno a arder} \textsuperscript{(10)}.

%%K t06-c4-03-1 p
3. --- Le supra del camino es ordinate; le \VUgras{lanterna}
\textsuperscript{(11)}, le \VUgras{cassetta de flammiferos}
\textsuperscript{(13)}, in le qual il ha le \VUgras{flammiferos}
\textsuperscript{(12)}, le \VUgras{candeliero} \textsuperscript{(14)} e le
\VUgras{portacandela} \textsuperscript{(15)} con un \VUgras{candela}
\textsuperscript{(16)} es a lor loco. Ma le grande \VUgras{situla pro carbon}
\textsuperscript{(18)}, plen de \VUgras{carbon de terra} \textsuperscript{(17)}
que le \VUgras{cocinera} \textsuperscript{(23)} justo plenava in le cellario,
restava in le medio del cocina.

%%K t06-c4-04-1 p
4. --- In ver, le povre femina debe hastar pro que le dejuner sia preste a
tempore. Ora illa es presso le \VUgras{furno de cocina} \textsuperscript{(19)} e
agita un \VUgras{sauce} \textsuperscript{(24)} que non debe arder troppo.

%%K t06-c4-05-1 p
5. --- In le \VUgras{furno} \textsuperscript{(20)} coce un torta, que illa
preparava pro le merenda del infantes. Super le furno de cocina pende le
\VUgras{coclear pro le potto} \textsuperscript{(21)} e le \VUgras{coclear
perforate} \textsuperscript{(22)}.

%%K t06-tit-8 sub
\VUsaut{4.0mm}
\VUpk{10.2pt}{40}{Le vasellamento.}
\VUsaut{3.0mm}

%%K t06-c4-06-1 p
6. --- Le \VUgras{batteria de cocina} \textsuperscript{(25)}, appendite in le
fundo del camera, es multo nette. Le belle \VUgras{cassarolas} de cupro
\textsuperscript{(26)}, le \VUgras{potto a lacte} \textsuperscript{(27)}, le
\VUgras{parve cribro} \textsuperscript{(28)} e le \VUgras{raspatorio}
\textsuperscript{(29)} esseva curosemente nettate.

%%K t06-c4-07-1 p
7. --- Le \VUgras{cassa a sal} \textsuperscript{(30)} es ponite sur le armario a
plattos presso le \VUgras{urceo a the} \textsuperscript{(31)}, e le \VUgras{grande
bottilia de oleo} \textsuperscript{(32)}. Le \VUgras{tiratorio}
\textsuperscript{(33)} probabilemente contine le utensiles de tabula e le
cultellos de cocina.

%%K t06-c4-08-1 p
8. --- Le \VUgras{scopa} \textsuperscript{(34)} e le \VUgras{plumaceo}
\textsuperscript{(35)} es in un angulo. Le cocinera justo usava le \VUgras{bloco
de ligno} \textsuperscript{(36)}, illa lo lassava presso le fenestra.

%%K t06-c4-09-1 p
9. --- Un \VUgras{toalia pro manos} \textsuperscript{(37)} pende al muro, presso
le \VUgras{imbuto} \textsuperscript{(38)}. In iste parte del cocina es
deponite le \VUgras{grillia} \textsuperscript{(39)}, le \VUgras{hacha}
\textsuperscript{(40)} e le \VUgras{planca a hachar} \textsuperscript{(41)}.
Postea, super le hacha, sur un \VUgras{tabuletta} \textsuperscript{(42)}, il ha
\VUgras{pottos} \textsuperscript{(43)}, le \VUgras{potto a bullir}
\textsuperscript{(44)} con su \VUgras{coperculo} \textsuperscript{(45)} e le
\VUgras{marmita} \textsuperscript{(46)}. Le \VUgras{patella}
\textsuperscript{(47)} pende presso.

%%K t06-c4-10-1 p
10. --- Solmente le \VUgras{robinetto} de cupro \textsuperscript{(48)} jecta un
brillo in iste angulo. Le \VUgras{lavamano} \textsuperscript{(49)} sur le qual
es ponite le grosse \VUgras{urceo} \textsuperscript{(50)}, es probabilemente
multo commode con su parve armario, in le qual on conserva le \VUgras{barril de
vinagre} \textsuperscript{(51)} providite de su \VUgras{robinetto}
\textsuperscript{(52)}, le \VUgras{cassa pro residuos} \textsuperscript{(53)} e
le \VUgras{plattos immunde} \textsuperscript{(54)}.

%%K t06-tit-9 sub
\VUsaut{4.0mm}
\VUpk{10.2pt}{40}{Altere objectos ponite in le cocina.}
\VUsaut{3.0mm}

%%K t06-c4-11-1 p
11. --- Le \VUgras{cuva} \textsuperscript{(55)} in le qual on lava le
\VUgras{legumines} \textsuperscript{(58)} e le \VUgras{caldiera}
\textsuperscript{(56)} de ferro fundite ponite sur le \VUgras{solo de tegulas}
\textsuperscript{(57)} probabilemente ha anque lor loco in iste armario.

%%K t06-c4-12-1 p
12. --- Le cocinera certemente non manca de furnos, proque ecce ancora, sur le
angulo del tabula, un \VUgras{furno a gas} \textsuperscript{(59)} sur le qual se
caleface aqua in un parve cassarola. Le \VUgras{vapor} \textsuperscript{(60)}
que escappa de illo nos indica, que probabilemente illo bulli. Verisimilemente
iste aqua essera usate pro le caffe, proque ecce sur le altere extremitate del
tabula, le \VUgras{cafetiera} \textsuperscript{(64)} e le \VUgras{molino a
caffe} \textsuperscript{(63)}.

%%K t06-c4-13-1 p
13. --- Le furno es connectite con le \VUgras{becco de gas}
\textsuperscript{(62)} per un longe \VUgras{tubetto de cauchu}
\textsuperscript{(61)}.

%%K t06-c4-14-1 p
14. --- Le \VUgras{servitrice a jornata} \textsuperscript{(66)} qui veni adjutar
le cocinera, prepara le legumines pro le cena. Quando illos essera mundate e
lavate, illa los ponera in le \VUgras{bassino} \textsuperscript{(65)}
attendente le hora pro cocer los.

%%K t06-tit-10 sub
\VUsaut{4.0mm}
{\centering\fontsize{10.2pt}{10.2pt}\selectfont\textls[40]{\scshape Le sala de banio.} \textit{(Camera de banio.)}\par}
\VUsaut{3.0mm}

%%K t06-c5-01-1 p
1. --- Nos ha a facer solmente un indiscretion pro cognoscer le principal
cameras del appartamento: vider le installation del sala de banio. Isto non
requirera un longe tempore, proque illo non es grande, e nos tosto sapera que
le \VUgras{cupa de banio} \textsuperscript{(1)} ha un bon \VUgras{calefactor de
banio} \textsuperscript{(2)}, que le \VUgras{concha pro sapon}
\textsuperscript{(3)} es attingibile per le mano del baniator e que le
\VUgras{portatoalias} \textsuperscript{(5)} certemente facera facile siccar le
\VUgras{toalias} \textsuperscript{(4)}.

%%K t06-c5-02-1 p
2. --- Iste camera ha su parietes coperite de \VUgras{faiencerias}
\textsuperscript{(6)}, que permitteva installar un \VUgras{apparato de ducha}
\textsuperscript{(7)} sin timer in alcun modo que le aqua, que resparge durante
su uso, damnificara le muros. Un \VUgras{parve bassino} \textsuperscript{(8)} de
zinc que servi pro le banios de pedes, completa le mobiliario.
\VUsaut{6.0mm}
\VUfilet{22mm}
